%%%%%%%%%%%%%%%%%%%%%%%%%%%%%%%%%%%%%%%%%%%%%%%%%%%%%%%%%%%%
%% CHAPTER 7 -- Why Bigger Sometimes Helps
%%%%%%%%%%%%%%%%%%%%%%%%%%%%%%%%%%%%%%%%%%%%%%%%%%%%%%%%%%%%

\chapter{Why Bigger Sometimes Helps}
\label{ntg:ch7}

\scenelabel
\begin{scenestrip}
\sceneitem{The problem}
\sceneitem{How we got here}
\sceneitem{Where we stand}
\sceneitem{What goes wrong}
\end{scenestrip}

\paragraph{The problem.} A modern frontier model has hundreds of billions of parameters,
sometimes more than a trillion. The internal arithmetic of such
a model does not fit on a single computer chip. Neither do the
intermediate values it produces while running, nor the
adjustments the training procedure makes to its parameters. A
single training run, as we noted in Chapter~2, takes weeks,
runs on tens of thousands of accelerator chips, and costs tens
to hundreds of millions of dollars. At inference time --- when
a user sends a query and waits for a response --- the same
model has to answer in seconds, often under load from thousands
of simultaneous queries.

The engineering that makes this possible is the subject of this
chapter. None of it is glamorous. Most of it is invisible to
the user. All of it determines whether a model trains in a
month or in a year, and whether the deployed system serves a
query at five cents or fifty. The same engineering, perversely,
is also what makes it economically rational to keep building
models larger and to keep finding cleverer ways to use them.
The answer to the question of how big a model should be is
inseparable from the answer to how efficiently the budget can
be spent.

\paragraph{How we got here.} The first models too large to fit on a single graphics card
were trained with a technique called \emph{data parallelism}.
The idea is simple: every worker holds a complete copy of the
model and processes a different slice of the training data, and
their gradient updates are averaged across workers. Data
parallelism is conceptually clean and runs out of room the
moment the model itself stops fitting on a single worker.

The first response was \emph{tensor parallelism}, popularised
by NVIDIA's Megatron-LM team in 2019. The idea is to split each
matrix multiplication across multiple workers, with the workers
exchanging partial results between them. The cost is the
communication between workers, which dominates how fast the
training can run. The benefit is that the model can be larger
than any single worker.

The next response was \emph{pipeline parallelism}, refined in
Google's GPipe work and used widely thereafter. The idea is to
put different layers of the model on different workers and to
run mini-batches in a staggered pipeline so that no worker sits
idle. The benefit is the same as tensor parallelism (the model
can be larger than a worker); the cost profile is different
(less communication-heavy, more sensitive to pipeline-bubble
inefficiencies).

A 2019 contribution from Microsoft Research called
\emph{ZeRO} --- the \emph{Zero Redundancy Optimizer} --- pointed
out that the bookkeeping the optimiser needs to do (the running
averages from Adam, for instance) was occupying more memory
than the model parameters themselves. ZeRO showed that this
bookkeeping could be split across workers without any loss of
correctness, and the libraries that implement the idea
(DeepSpeed, PyTorch FSDP) became standard tools.

The list goes on. There is a parallelism for splitting the
sequence across workers (when the context window is the
binding constraint). There is a parallelism for splitting the
specialist subnetworks of a Mixture-of-Experts model (which we
discuss below). There are libraries that overlap communication
between workers with the computation that does not depend on
that communication. There are inference-time tricks that batch
queries together to keep the graphics cards saturated. None of
this was in the original Transformer paper. All of it is the
work of the engineering community in the years since. The
infrastructure surrounding a modern model is, by some measures,
an order of magnitude more complex than the model it serves.

The other half of the chapter is sparsity. The idea that a
neural network could be made of specialist subnetworks, with
only a fraction of them active for any given input, is older
than the deep learning era. Robert Jacobs, Michael Jordan, and
their collaborators introduced \emph{adaptive mixtures of
experts} in 1991. The idea was to train a small number of
specialist networks alongside a learned ``gate'' that decided
which specialist would handle each input. The idea sat largely
dormant during the era of dense networks. Noam Shazeer and
collaborators at Google revived it in 2017 with the
\emph{Sparsely-Gated Mixture-of-Experts Layer}, which scaled a
language model to 137 billion parameters with only a small
fraction of those parameters active for any given token. The
\emph{Switch Transformer} in 2021 simplified the routing
further. Mistral AI's \emph{Mixtral 8x7B} in December 2023 brought
mixture-of-experts models into wide open-weight use. The
appeal is clean: you get the parameter count of a much larger
dense model at the per-token compute of a much smaller one,
because most of the experts sit out most of the time.

\paragraph{Where we stand.} A modern training stack composes several forms of parallelism
at once, with the choice tuned to the specific cluster's
network topology. Mixed-precision arithmetic --- using lower-
precision numbers for most calculations and higher-precision
numbers only where stability requires --- reduces memory and
increases throughput. Mixture-of-Experts layers, in production
designs, route each token to a small fixed number of experts
(typically two out of eight or sixteen) with auxiliary losses
that discourage the routing from collapsing onto a few
favourites. Inference serving stacks --- vLLM, TensorRT-LLM,
SGLang, and their relatives --- batch queries together,
schedule them with continuous batching, and decode with
techniques like speculative or parallel sampling.

The combined effect is that a frontier model in 2024 trains for
roughly the same number of weeks on a cluster of roughly the
same size as a frontier model in 2020, while being substantially
more capable. The gain is partly from larger and better-curated
data, partly from architectural improvements, partly from
algorithmic improvements in the training procedure, and partly
from the engineering work described above. Drawing precise
attribution between these factors is difficult; the qualitative
point is that the field is making progress on several axes at
once, and that ``how much compute did this cost'' is a less
useful summary than it used to be.

\paragraph{What goes wrong.} The dramatic failure modes of distributed training are visible:
a network partition that hangs a job for hours before anyone
notices, a subtle ordering bug in the communication between
workers that corrupts gradients silently and produces a model
that trains badly, a mixed-precision calculation that underflows
or overflows and produces a run that has to be restarted from
the last checkpoint. These failures are part of the operational
reality of running a frontier training cluster, and the teams
that operate such clusters spend a non-trivial fraction of
their time on monitoring, recovery, and debugging.

Mixture-of-experts models bring their own failure modes. The
routing can collapse onto a handful of experts, leaving the
rest of the experts under-trained and the model effectively
smaller than its parameter count suggests. The routing can
collapse the other way, sending too many tokens to a single
expert and forcing some of them to be dropped. The auxiliary
loss that prevents these collapses is itself a hyperparameter
whose tuning is delicate. None of this is unfixable, but all of
it adds operational risk that dense models do not carry.

The deeper problem is that the scaling laws of Chapter~2 are a
description of how the loss decreases with resources, not a
guarantee that the deployed model will be useful. A model that
has trained to a lower loss is, on average, more capable than a
model that has not, but the relationship between the loss and
any specific deployment property is loose. We will see this
again in Chapter~10 when we discuss evaluation, and in
Chapter~13 when we discuss the gap between what the model
optimises and what the user wants.

\section{The idea, in plain language}

\subsection*{The Chinchilla insight, revisited}

Chapter~2 introduced the Chinchilla scaling law in passing. It
is worth stating it more carefully here because it is the
single most consequential empirical result in the field's
recent history.

Before the Chinchilla paper of 2022, the dominant view in the
field --- following the Kaplan paper of 2020 --- was that, given
a fixed compute budget, you should spend most of it on
parameters and relatively less on training data. The largest
models of 2020 and 2021 followed this prescription. GPT-3, at
175 billion parameters, was trained on about 300 billion
tokens.

The Chinchilla paper, by Jordan Hoffmann and a team at
DeepMind, ran a more careful version of the experiment and
reached a different conclusion. They found that, given a fixed
compute budget, you should scale parameters and tokens roughly
together, not parameters faster than tokens. Specifically: for
each parameter in your model, you should train on roughly
twenty tokens. By this measure, GPT-3 had been substantially
under-trained. A model of about 70 billion parameters trained
on 1.4 trillion tokens (the Chinchilla model itself) would, on
the same compute budget, outperform GPT-3 substantially. The
prediction was tested. It held.

The Chinchilla finding changed how the industry budgets compute.
Most modern frontier models are smaller in parameter count than
GPT-3 and trained on much more data. Llama 3, released by Meta
in 2024, was trained on more than fifteen trillion tokens; this
is roughly fifty times the GPT-3 training set, on a model that
has fewer parameters than GPT-3 had. The combination produces a
much more capable model. The qualitative point is that, in the
post-Chinchilla world, it does not pay to brag about parameter
count. The relevant numbers are how much data was used, of what
quality, and what compute budget was applied to it.

\subsection*{Mixture-of-experts as a design pattern}

The mixture-of-experts idea is best understood as the answer to
a specific economic question: how can you have a model with
many parameters (which seems to help capability) without paying
the inference cost of activating all of those parameters on
every query?

The answer is to organise some of the model's parameters into
specialist subnetworks --- the ``experts'' --- and to put a
small \emph{router} network in front of them whose only job is
to decide, for each new token, which experts should be
activated. The total number of parameters in the model can be
very large; the number of parameters actually used in producing
any single token's output is much smaller, because only a few
experts are activated at a time.

The design has appealing properties. The total number of
parameters scales with the number of experts; the per-token
compute scales with the number of experts active per token,
which is much smaller. A model with eight experts, of which two
are activated per token, has roughly the parameter count of an
eight-expert model and the per-token compute of a two-expert
one. This is a real efficiency gain.

The design has corresponding costs. All the experts have to be
held in memory at inference time, even though most of them are
inactive on most tokens. So the memory cost of an MoE model
scales with the parameter count, even though the compute cost
scales with the active fraction. For deployments where memory
is the binding constraint --- consumer hardware, mobile devices,
cost-sensitive serving --- this is a serious limitation. For
deployments where compute is the binding constraint --- large-
scale serving on dedicated hardware --- the trade-off is
favourable.

The other cost is operational. The routing has to be carefully
trained to avoid collapse, the auxiliary loss has to be tuned,
and the inference engine has to handle the variable routing
efficiently. These are all solved problems in the open
literature, but they add complexity that a dense model does
not have.

\subsection*{What scaling does not buy you}

It is easy, when reading announcements about new frontier
models, to fall into the assumption that capability is a single
dimension and that bigger is uniformly better. Both halves of
the assumption are wrong.

Capability is a multidimensional thing. A model can be better
at translation and worse at code; better at conversation and
worse at mathematical reasoning; better at English and worse
at Hindi. The relative balance among these capabilities depends
on the training data mix, on the architecture, on the post-
training, and on a great many other choices. ``Better than X''
is a useful summary only when the ``better at what'' is also
specified.

Bigger is uniformly better only along the axes that scale with
the chosen training data. If the training corpus does not
contain much Mongolian, no amount of additional parameters will
turn the model into a competent Mongolian-language assistant.
If the training corpus does not contain much recent biological
research, no amount of scale will let the model reason about
post-training-cut-off discoveries. The scaling laws describe
how the loss falls with resources, on the distribution the
model is trained on. They do not promise behaviour outside that
distribution. We will return to this in Chapter~10 when we
discuss retrieval as a way of importing information the model
was not trained on.

Bigger is also not free. The energy and water cost of frontier
training has become a public concern; the absolute numbers are
disputed, but the trajectory is unambiguous. The economic cost
of training the next generation of models is, to a degree that
the industry does not always discuss publicly, restricting the
field of organisations that can produce frontier models to a
handful. Both of these are second-order consequences of
scaling, and both are likely to shape the regulatory and
geopolitical landscape over the next decade.

\section{The engineering consequence}

\paragraph{The cost of inference is dominated by serving, not
training.} A frontier model that cost a hundred million
dollars to train will be served to billions of queries over
its life. The training cost is amortised. The serving cost
per query --- the cost of running the model once --- is what
determines the unit economics of any product built on top of
it. For most product teams the relevant number is not ``how
much did the training cost'' but ``how much does this query
cost'', and the latter is mostly a function of the
infrastructure described in this chapter.

\paragraph{Mixture-of-experts is an attractive design for some
deployments.} Where the deployment can afford the memory
footprint of all experts but is constrained on per-query
compute, MoE provides better quality per dollar than the dense
alternative. Where the deployment is constrained on memory --- on
device, on consumer hardware, on cost-sensitive serving --- a
dense model of the same active-parameter count is usually
preferable.

\paragraph{Vendor scaling claims need translation.} A model
that is described as ``trillion-parameter'' may be a dense
trillion-parameter model (which would be the largest ever
publicly disclosed) or it may be a mixture-of-experts model
whose total parameter count includes inactive experts. The two
have different implications for capability, cost, and the
operational complexity the vendor will pass through to the
customer. A non-technical reader evaluating such claims should
ask the vendor to disaggregate the parameter count by
``parameters total'' and ``parameters active per token''.

\paragraph{Scaling-law extrapolation is risky beyond the
measured range.} The empirical regularity that loss falls
predictably with resources has held over many orders of
magnitude. It need not hold over the next several. There are
plausible reasons --- exhaustion of high-quality training
data, diminishing returns from each doubling, fundamental
limits of next-token prediction as a training objective ---
why the curve might bend. There are also plausible reasons
why it might not. The honest position is that the field's
ability to predict the capability of the next generation from
the current one is mediocre, and that anyone making firm
predictions in this area has overweighted whichever piece of
evidence supports their position.

\section{What goes wrong, and why}

\paragraph{The data wall.} The Chinchilla recipe says that
making models better requires more training data, in roughly
proportional measure to the number of parameters. The supply
of high-quality training data on the public internet is finite.
Estimates of when the supply will be exhausted vary, but
several recent analyses suggest that the most ambitious
training runs of the next few years will require either
licensed proprietary data, synthetic data generated by other
models, or both. Each of these has its own complications.
Licensed data is expensive and concentrates value in whoever
holds the rights. Synthetic data risks producing a closed loop
in which models are trained on the output of other models, with
unknown long-term consequences. The data wall is not a single
event; it is a slow tightening of the constraint.

\paragraph{The energy and infrastructure constraint.} A
frontier training run consumes electricity on the scale of a
small town for the duration of the run. The data centres that
host these runs are concentrated in regions with favourable
electricity prices and grid capacity, and several recent
expansions have met local resistance over water, land, and
power use. These are not theoretical concerns. They appear in
the permit records of data-centre proposals and in the
financial disclosures of the companies building them. The
direction of the field is increasingly shaped by the
availability of its physical infrastructure.

\paragraph{Sparse-routing collapse.} Mixture-of-experts models
can degrade in subtle ways if the routing is not carefully
managed. An expert that has stopped receiving tokens has
stopped learning; an expert that is receiving more tokens than
it can handle has stopped being a specialist. These failures
are not visible in the loss curve directly --- the loss may
still decrease --- but they are visible in the model's
deployment behaviour as uneven capability across topics. An
operator who knows the model is MoE-based should ask whether
the developer has audited expert utilisation post-training.

\paragraph{Geopolitical concentration.} The hardware required
to train a frontier model is, at the time of writing, almost
exclusively produced by NVIDIA and almost exclusively
fabricated by TSMC in Taiwan. The political ramifications of
this concentration are out of scope for this guide, but the
operational implications are not: an organisation whose product
depends on a frontier model is, indirectly, depending on a
hardware supply chain whose continuity is not guaranteed. This
is one of the under-discussed risks in the LLM-product space,
and a non-technical reader evaluating long-term commitments to
LLM-based products should be aware of it.

\section*{Bridges}
\addcontentsline{toc}{section}{Bridges}

This chapter built on Chapter~2 (the scaling-law framing) and
Chapter~6 (the pretraining recipe) and described the
engineering by which the scaling is actually realised. It also
described the mixture-of-experts design as one of the more
consequential recent moves to make scaling more efficient.
Chapter~8 looks at the other half of the cost-control story:
how to make a trained model adaptable without retraining it
from scratch. Chapter~10 looks at how to make a trained model
useful on data it was not trained on. Chapter~11 looks at the
governance implications of the geopolitical and infrastructural
concentration described above. The next chapter --- \emph{Adapting
a Giant Model} --- explains how a model trained at industrial
scale ends up answering questions about your specific company,
your specific dataset, your specific tone.
